\chapter{The Financial System, Institutions, Markets, and Regulation}

\section{What the system does}
The financial system connects savers, borrowers, risk bearers, and payment users. Its core functions are to allocate capital, transfer risk, provide liquidity, produce information through prices, and settle obligations. These functions are performed by institutions and market infrastructures rather than by a single market.

\begin{itemize}
  \item \textbf{Intermediaries} such as banks, insurers, pension funds, and asset managers pool funds or transform risks.
  \item \textbf{Markets} allow claims to be issued and traded. Primary markets finance issuers; secondary markets let investors change positions.
  \item \textbf{Infrastructures} include payment systems, securities depositories, clearing houses, exchanges, and trade repositories.
  \item \textbf{Authorities} set rules, supervise institutions, resolve failures, protect certain users, and sometimes provide emergency liquidity.
\end{itemize}

The categories overlap. A bank can be an intermediary, a market participant, a payment provider, and a regulated entity. A central counterparty reduces bilateral counterparty exposure by becoming the buyer to every seller and the seller to every buyer, but it concentrates other risks that must be managed with margin, default funds, and governance.

\section{Markets and instruments}
\keyterm{Money markets} concern short maturities and instruments such as bills, commercial paper, certificates of deposit, and repurchase agreements. \keyterm{Capital markets} support longer-term debt and equity financing. \keyterm{Foreign-exchange markets} exchange currencies and settle international trade and investment. \keyterm{Commodity markets} transfer price and delivery risks for physical goods. \keyterm{Derivative markets} trade contracts whose value depends on an underlying asset, rate, index, or event.

Market design affects behaviour. An exchange usually standardises contracts, publishes quotes, applies eligibility rules, and uses a matching mechanism. Over-the-counter trading permits customised terms but creates bilateral exposure, valuation disputes, and less transparent prices. Electronic trading can reduce some frictions while introducing speed, connectivity, and operational risks.

\section{Banks and money creation}
A bank's balance sheet contains assets such as loans and securities, liabilities such as deposits and wholesale borrowing, and equity capital. When a bank makes a loan, it commonly records both a loan asset and a deposit liability; the transaction is not simply a transfer of pre-existing cash. Settlement and liquidity constraints still matter, and broad money is linked to the banking system and central-bank operations rather than to an unlimited ability to lend.

The maturity transformation of short-term liabilities into longer-term assets can support investment, but it creates run risk. A depositor may demand cash before a loan is repaid. Liquidity regulation, capital requirements, deposit insurance, supervision, and lender-of-last-resort arrangements address parts of this problem; none eliminates the possibility of failure.

\section{Funds, insurers, and market makers}
An investment fund pools investor money under a mandate. A mutual fund typically issues and redeems units at a net asset value, while an exchange-traded fund trades on an exchange and usually relies on authorised participants to create or redeem units. Fees, tracking error, liquidity, tax treatment, and the legal structure matter as much as the name of the product.

An insurer collects premiums and pays uncertain claims. It diversifies risks across policyholders but must manage underwriting, reserving, investment, catastrophe, and liquidity risks. A market maker provides two-sided quotes and earns compensation for inventory, adverse-selection, funding, and operational risks. The quoted spread is not pure profit.

\section{Regulation: aims and limits}
Regulation pursues several objectives: investor protection, market integrity, safety and soundness, financial stability, fair disclosure, competition, and the prevention of crime. Different authorities divide these goals differently. A single transaction can therefore be affected by securities law, banking rules, derivatives rules, accounting standards, consumer law, data protection, tax law, and insolvency law.

The Basel Committee describes Basel III as an internationally agreed set of measures intended to strengthen bank regulation, supervision, and risk management after the global financial crisis \cite{basel3}. The Basel Framework contains prudential standards but is implemented through national or regional law. It should not be treated as a single worldwide statute.

The European Central Bank describes financial stability in terms of a system of intermediaries, markets, and infrastructures that can withstand shocks without severe disruption to real activity \cite{ecb2026}. Its Financial Stability Review is a dated assessment of euro-area vulnerabilities, not a prediction of market returns. In the United States, the SEC, CFTC, Federal Reserve, Treasury, banking supervisors, state authorities, and self-regulatory organisations have distinct roles. The legal answer for a real transaction depends on jurisdiction, product, counterparty, and facts.

\section{Settlement and custody}
A trade is not complete when a screen shows an execution. The parties must confirm the contract, calculate obligations, exchange collateral where required, settle cash and securities, record ownership, and handle corporate actions. \keyterm{Settlement risk} is the possibility that one side performs while the other does not. Delivery-versus-payment reduces this risk by linking delivery of a security to payment.

\section{Information and incentives}
Financial institutions solve information problems but can also create them. Lenders screen and monitor borrowers. Exchanges publish prices. Rating agencies analyse credit. Auditors provide assurance over financial statements. Yet each agent has incentives, limited information, and potential conflicts. A regulation that improves one objective may create a trade-off elsewhere; for example, higher capital can improve resilience but may change intermediation costs and market structure.

\source{Institutional descriptions are consistent with SEC investor education, TreasuryDirect documentation, the BIS Basel Framework, IOSCO principles for financial market infrastructures, and the ECB's financial-stability materials. Jurisdiction-specific rules are intentionally not reduced to investment advice.}

\begin{exerciseblock}
\begin{enumerate}
  \item Distinguish a primary market from a secondary market using a new share issue as an example.
  \item Why can maturity transformation create value and fragility at the same time?
  \item Name two risks that a central counterparty reduces and two risks it can concentrate.
  \item Explain why a regulation described as ``international'' may still require local implementing law.
\end{enumerate}
\end{exerciseblock}

